\documentclass[tikz]{standalone}
\input{../typography.tex.preamble}
\input{../tikz.tex.preamble}
\usepackage{cancel}
\usetikzlibrary{decorations.pathreplacing}

\tikzset{
  declare function={
    f(\x) = \x^3/12 - \x^2/6 - \x + 1;
    a = -2/3*(sqrt(10) - 1);
    b = +2/3*(sqrt(10) + 1);
    c = 4;
  }
}

\pgfplotsset{
  extrema/.style={
    axis lines=middle,
    axis equal image,
    enlargelimits=true,
    xlabel={\(x\)}, xlabel style={anchor=west},
    ylabel={\(y\)}, ylabel style={anchor=south},
    label style={at={(ticklabel* cs:1)}, font={\tiny}},
    xtick=\empty,
    ytick=\empty,
    ticklabel style={font=\tiny},
    grid=none, grid style={gray!20},
    smooth, samples=100, no markers,
  }
}

\newcommand{\basecurve}{%
  \addplot[thick, domain=-4:5] {f(x)};
  \addplot[thick] coordinates {(5, {f(5)}) (7, {f(5)})};
  \draw[fill=white] (7, {f(5)}) circle[radius=2pt];
  \draw[fill=black] (7, 3) circle[radius=2pt];
  \addplot[thick] coordinates {(7, 3) (9, 4)};
  \addplot[thick] coordinates {(9, 4) (10, 2)};
  \draw[fill=black] (10, 2) circle[radius=2pt];%
}

\begin{document}

% extrema_example_blank
\begin{tikzpicture}
  \begin{axis}[extrema]
    \basecurve
  \end{axis}
\end{tikzpicture}

% extrema_example_1_1
\begin{tikzpicture}
  \begin{axis}[extrema]
    \draw[main, dotted, fill=attn!10] (-4, 4) rectangle (6, 3);
    \node at (axis cs:1,3.5) {\sffamily\footnotesize \phantom{because} \(f'(c) = 0\)};
    \node[below] (C) at (axis cs:a,1) {\sffamily\footnotesize \(c\)};
    \draw[attn, ->] (axis cs:1,3) -- (C);
    \basecurve
    % smooth local max
    \addplot[main] coordinates {(a,0.2) (a,-0.2)};
    \fill[supp] ({a}, {f(a)}) circle[radius=0.2];
    \addplot[supp, dashed] coordinates { (a-2, {f(a)}) ({a+2}, {f(a)}) };
  \end{axis}
\end{tikzpicture}

% extrema_example_1
\begin{tikzpicture}
  \begin{axis}[extrema]
    \draw[main, dotted, fill=attn!10] (-4, 4) rectangle (6, 3);
    \node at (axis cs:1,3.5) {\sffamily\footnotesize because \(f'(c) = 0\)};
    \node[below] (C) at (axis cs:a,1) {\sffamily\footnotesize \(c\)};
    \draw[attn, ->] (axis cs:1,3) -- (C);
    \basecurve
    % smooth local max
    \addplot[main] coordinates {(a,0.2) (a,-0.2)};
    \node[main, below] at (axis cs:a,-0.2) {\tiny\sffamily CN};
    \fill[supp] ({a}, {f(a)}) circle[radius=0.2];
    \addplot[supp, dashed] coordinates { (a-2, {f(a)}) ({a+2}, {f(a)}) };
  \end{axis}
\end{tikzpicture}

% extrema_example_2_1
\begin{tikzpicture}
  \begin{axis}[extrema]
    \draw[main, dotted, fill=attn!10] (-4, 4) rectangle (6, 3);
    \node at (axis cs:1,3.5) {\sffamily\footnotesize \phantom{because} \(f'(c) = 0\)};
    \node[above] (C) at (axis cs:b,-1) {\sffamily\footnotesize \(c\)};
    \draw[attn, ->] (axis cs:1,3) -- (C);
    \basecurve
    % smooth local max
    \addplot[main] coordinates {(a,0.2) (a,-0.2)};
    \node[main, below] at (axis cs:a,-0.2) {\tiny\sffamily CN};
    % smooth local min
    \addplot[main] coordinates {(b,-0.2) (b,0.2)};
    \fill[supp] (b, {f(b)}) circle[radius=0.2];
    \addplot[supp, dashed] coordinates { (b-2, {f(b)}) ({b+2}, {f(b)}) };
  \end{axis}
\end{tikzpicture}

% extrema_example_2
\begin{tikzpicture}
  \begin{axis}[extrema]
    \draw[main, dotted, fill=attn!10] (-4, 4) rectangle (6, 3);
    \node at (axis cs:1,3.5) {\sffamily\footnotesize because \(f'(c) = 0\)};
    \node[above] (C) at (axis cs:b,-1) {\sffamily\footnotesize \(c\)};
    \draw[attn, ->] (axis cs:1,3) -- (C);
    \basecurve
    % smooth local max
    \addplot[main] coordinates {(a,0.2) (a,-0.2)};
    \node[main, below] at (axis cs:a,-0.2) {\tiny\sffamily CN};
    % smooth local min
    \addplot[main] coordinates {(b,-0.2) (b,0.2)};
    \node[main, above] at (axis cs:b,0.2) {\tiny\sffamily CN};
    \fill[supp] (b, {f(b)}) circle[radius=0.2];
    \addplot[supp, dashed] coordinates { (b-2, {f(b)}) ({b+2}, {f(b)}) };
  \end{axis}
\end{tikzpicture}

% extrema_example_3_1
\begin{tikzpicture}
  \begin{axis}[extrema]
    \draw[main, dotted, fill=attn!10] (-4, 4) rectangle (6, 3);
    \node at (axis cs:1,3.5) {\sffamily\footnotesize \phantom{because} \(f'(c)\) does not exist};
    \node[above] (C) at (axis cs:5,-1) {\sffamily\footnotesize \(c\)};
    \draw[attn, ->] (axis cs:1,3) -- (C);
    \basecurve
    % smooth local max
    \addplot[main] coordinates {(a,0.2) (a,-0.2)};
    \node[main, below] at (axis cs:a,-0.2) {\tiny\sffamily CN};
    % smooth local min
    \addplot[main] coordinates {(b,-0.2) (b,0.2)};
    \node[main, above] at (axis cs:b,0.2) {\tiny\sffamily CN};
    % flat
    \addplot[main] coordinates {(5,0.2) (5,-0.2)};
  \end{axis}
\end{tikzpicture}

% extrema_example_3
\begin{tikzpicture}
  \begin{axis}[extrema]
    \draw[main, dotted, fill=attn!10] (-4, 4) rectangle (6, 3);
    \node at (axis cs:1,3.5) {\sffamily\footnotesize because \(f'(c)\) does not exist};
    \node[above] (C) at (axis cs:5,-1) {\sffamily\footnotesize \(c\)};
    \draw[attn, ->] (axis cs:1,3) -- (C);
    \basecurve
    % smooth local max
    \addplot[main] coordinates {(a,0.2) (a,-0.2)};
    \node[main, below] at (axis cs:a,-0.2) {\tiny\sffamily CN};
    % smooth local min
    \addplot[main] coordinates {(b,-0.2) (b,0.2)};
    \node[main, above] at (axis cs:b,0.2) {\tiny\sffamily CN};
    % flat
    \addplot[main] coordinates {(5,0.2) (5,-0.2)};
    \node[main, above] at (axis cs:5,0) {\tiny\sffamily CN};
  \end{axis}
\end{tikzpicture}

% extrema_example_4_1
\begin{tikzpicture}
  \begin{axis}[extrema]
    \draw[main, dotted, fill=attn!10] (-4, 4) rectangle (6, 3);
    \node at (axis cs:1,3.5) {\sffamily\footnotesize \phantom{because} \(f'(c)\) does not exist};
    \node[above] (C) at (axis cs:7,-1) {\sffamily\footnotesize \(c\)};
    \draw[attn, ->] (axis cs:1,3) -- (C);
    \basecurve
    % smooth local max
    \addplot[main] coordinates {(a,0.2) (a,-0.2)};
    \node[main, below] at (axis cs:a,-0.2) {\tiny\sffamily CN};
    % smooth local min
    \addplot[main] coordinates {(b,-0.2) (b,0.2)};
    \node[main, above] at (axis cs:b,0.2) {\tiny\sffamily CN};
    % flat
    \addplot[main] coordinates {(5,0.2) (5,-0.2)};
    \addplot[main] coordinates {(7,0.2) (7,-0.2)};
    \node[main, above] at (axis cs:5,0) {\tiny\sffamily CN};
  \end{axis}
\end{tikzpicture}

% extrema_example_4
\begin{tikzpicture}
  \begin{axis}[extrema]
    \draw[main, dotted, fill=attn!10] (-4, 4) rectangle (6, 3);
    \node at (axis cs:1,3.5) {\sffamily\footnotesize because \(f'(c)\) does not exist};
    \node[above] (C) at (axis cs:7,-1) {\sffamily\footnotesize \(c\)};
    \draw[attn, ->] (axis cs:1,3) -- (C);
    \basecurve
    % smooth local max
    \addplot[main] coordinates {(a,0.2) (a,-0.2)};
    \node[main, below] at (axis cs:a,-0.2) {\tiny\sffamily CN};
    % smooth local min
    \addplot[main] coordinates {(b,-0.2) (b,0.2)};
    \node[main, above] at (axis cs:b,0.2) {\tiny\sffamily CN};
    % flat
    \addplot[main] coordinates {(5,0.2) (5,-0.2)};
    \addplot[main] coordinates {(7,0.2) (7,-0.2)};
    \node[main, above] at (axis cs:5,0) {\tiny\sffamily CN};
    \node[main, above] at (axis cs:7,0) {\tiny\sffamily CN};
  \end{axis}
\end{tikzpicture}

% extrema_example_5_1
\begin{tikzpicture}
  \begin{axis}[extrema]
    \draw[main, dotted, fill=attn!10] (-4, 4) rectangle (6, 3);
    \node at (axis cs:1,3.5) {\sffamily\footnotesize \phantom{because} \(f'(c)\) does not exist};
    \node[below] (C) at (axis cs:9,1) {\sffamily\footnotesize \(c\)};
    \draw[attn, ->] (axis cs:1,3) -- (C);
    \basecurve
    % smooth local max
    \addplot[main] coordinates {(a,0.2) (a,-0.2)};
    \node[main, below] at (axis cs:a,-0.2) {\tiny\sffamily CN};
    % smooth local min
    \addplot[main] coordinates {(b,-0.2) (b,0.2)};
    \node[main, above] at (axis cs:b,0.2) {\tiny\sffamily CN};
    % flat
    \addplot[main] coordinates {(5,0.2) (5,-0.2)};
    \addplot[main] coordinates {(7,0.2) (7,-0.2)};
    \node[main, above] at (axis cs:5,0) {\tiny\sffamily CN};
    \node[main, above] at (axis cs:7,0) {\tiny\sffamily CN};
    % absolute max
    \addplot[main] coordinates {(9,0.2) (9,-0.2)};
  \end{axis}
\end{tikzpicture}

% extrema_example_5
\begin{tikzpicture}
  \begin{axis}[extrema]
    \draw[main, dotted, fill=attn!10] (-4, 4) rectangle (6, 3);
    \node at (axis cs:1,3.5) {\sffamily\footnotesize because \(f'(c)\) does not exist};
    \node[below] (C) at (axis cs:9,1) {\sffamily\footnotesize \(c\)};
    \draw[attn, ->] (axis cs:1,3) -- (C);
    \basecurve
    % smooth local max
    \addplot[main] coordinates {(a,0.2) (a,-0.2)};
    \node[main, below] at (axis cs:a,-0.2) {\tiny\sffamily CN};
    % smooth local min
    \addplot[main] coordinates {(b,-0.2) (b,0.2)};
    \node[main, above] at (axis cs:b,0.2) {\tiny\sffamily CN};
    % flat
    \addplot[main] coordinates {(5,0.2) (5,-0.2)};
    \addplot[main] coordinates {(7,0.2) (7,-0.2)};
    \node[main, above] at (axis cs:5,0) {\tiny\sffamily CN};
    \node[main, above] at (axis cs:7,0) {\tiny\sffamily CN};
    % absolute max
    \addplot[main] coordinates {(9,0.2) (9,-0.2)};
    \node[main, below] at (axis cs:9,-0.2) {\tiny\sffamily CN};
  \end{axis}
\end{tikzpicture}

% extrema_example_6_1
\begin{tikzpicture}
  \begin{axis}[extrema]
    \draw[main, dotted, fill=attn!10] (-4, 4) rectangle (6, 3);
    \node at (axis cs:1,3.5) {\sffamily\footnotesize notice \(f'(c)\) does not exist};
    \node[above] (C) at (axis cs:10,-1) {\sffamily\footnotesize \(c\)};
    \draw[attn, ->] (axis cs:1,3) -- (C);
    \basecurve
    % smooth local max
    \addplot[main] coordinates {(a,0.2) (a,-0.2)};
    \node[main, below] at (axis cs:a,-0.2) {\tiny\sffamily CN};
    % smooth local min
    \addplot[main] coordinates {(b,-0.2) (b,0.2)};
    \node[main, above] at (axis cs:b,0.2) {\tiny\sffamily CN};
    % flat
    \addplot[main] coordinates {(5,0.2) (5,-0.2)};
    \addplot[main] coordinates {(7,0.2) (7,-0.2)};
    \node[main, above] at (axis cs:5,0) {\tiny\sffamily CN};
    \node[main, above] at (axis cs:7,0) {\tiny\sffamily CN};
    % absolute max
    \addplot[main] coordinates {(9,0.2) (9,-0.2)};
    \node[main, below] at (axis cs:9,-0.2) {\tiny\sffamily CN};
    \addplot[warn] coordinates {(10,0.2) (10,-0.2)};
    \node[warn, above] at (axis cs:10,0.2) {\tiny\sffamily CN??};
  \end{axis}
\end{tikzpicture}

% extrema_example_6
\begin{tikzpicture}
  \begin{axis}[extrema]
    \draw[main, dotted, fill=attn!10] (-4, 4) rectangle (6, 3);
    \node at (axis cs:1,3.5) {\sffamily\footnotesize \hlwarn{CN cannot be an endpoint}};
    \node[above] (C) at (axis cs:10,-1) {\sffamily\footnotesize \(c\)};
    \draw[attn, ->] (axis cs:1,3) -- (C);
    \basecurve
    % smooth local max
    \addplot[main] coordinates {(a,0.2) (a,-0.2)};
    \node[main, below] at (axis cs:a,-0.2) {\tiny\sffamily CN};
    % smooth local min
    \addplot[main] coordinates {(b,-0.2) (b,0.2)};
    \node[main, above] at (axis cs:b,0.2) {\tiny\sffamily CN};
    % flat
    \addplot[main] coordinates {(5,0.2) (5,-0.2)};
    \addplot[main] coordinates {(7,0.2) (7,-0.2)};
    \node[main, above] at (axis cs:5,0) {\tiny\sffamily CN};
    \node[main, above] at (axis cs:7,0) {\tiny\sffamily CN};
    % absolute max
    \addplot[main] coordinates {(9,0.2) (9,-0.2)};
    \node[main, below] at (axis cs:9,-0.2) {\tiny\sffamily CN};
    \addplot[warn] coordinates {(10,0.2) (10,-0.2)};
    \node[warn, above] at (axis cs:10,0.2) {\tiny\sffamily \xcancel{CN}};
  \end{axis}
\end{tikzpicture}

% extrema_example_7_1
\begin{tikzpicture}
  \begin{axis}[extrema]
    \basecurve
    % smooth local max
    \addplot[main] coordinates {(a,0.2) (a,-0.2)};
    \node[main, below] at (axis cs:a,-0.2) {\tiny\sffamily CN};
    % smooth local min
    \addplot[main] coordinates {(b,-0.2) (b,0.2)};
    \node[main, above] at (axis cs:b,0.2) {\tiny\sffamily CN};
    % flat
    \addplot[main] coordinates {(5,0.2) (5,-0.2)};
    \addplot[main] coordinates {(7,0.2) (7,-0.2)};
    \node[main, above] at (axis cs:5,0) {\tiny\sffamily CN};
    \node[main, above] at (axis cs:7,0) {\tiny\sffamily CN};
    % absolute max
    \addplot[main] coordinates {(9,0.2) (9,-0.2)};
    \node[main, below] at (axis cs:9,-0.2) {\tiny\sffamily CN};
    \addplot[warn] coordinates {(10,0.2) (10,-0.2)};
    \node[warn, above] at (axis cs:10,0.2) {\tiny\sffamily \xcancel{CN}};
  \end{axis}
\end{tikzpicture}

% extrema_example_7
\begin{tikzpicture}
  \begin{axis}[extrema]
    \basecurve
    % smooth local max
    \addplot[main] coordinates {(a,0.2) (a,-0.2)};
    \node[main, below] at (axis cs:a,-0.2) {\tiny\sffamily CN};
    % smooth local min
    \addplot[main] coordinates {(b,-0.2) (b,0.2)};
    \node[main, above] at (axis cs:b,0.2) {\tiny\sffamily CN};
    % flat
    \addplot[main] coordinates {(5,0.2) (5,-0.2)};
    \addplot[main] coordinates {(7,0.2) (7,-0.2)};
    \node[main, above] at (axis cs:5,0) {\tiny\sffamily CN};
    \node[main, above] at (axis cs:7,0) {\tiny\sffamily CN};
    \draw[main, ultra thick] (5,0) -- (7,0);
    \draw[main, decorate, decoration={brace, raise=0.2cm, mirror}] (5,0) -- (7,0);
    \node[main, below] at (axis cs:6,-0.6) {\tiny\sffamily all CNs};
    % absolute max
    \addplot[main] coordinates {(9,0.2) (9,-0.2)};
    \node[main, below] at (axis cs:9,-0.2) {\tiny\sffamily CN};
    \addplot[warn] coordinates {(10,0.2) (10,-0.2)};
    \node[warn, above] at (axis cs:10,0.2) {\tiny\sffamily \xcancel{CN}};
  \end{axis}
\end{tikzpicture}

% extrema_example
\begin{tikzpicture}
  \begin{axis}[extrema]
    \basecurve
    % smooth local max
    \addplot[main] coordinates {(a,0.2) (a,-0.2)};
    \node[main, below] at (axis cs:a,-0.2) {\tiny\sffamily CN};
    % smooth local min
    \addplot[main] coordinates {(b,-0.2) (b,0.2)};
    \node[main, above] at (axis cs:b,0.2) {\tiny\sffamily CN};
    % flat
    \addplot[main] coordinates {(5,0.2) (5,-0.2)};
    \addplot[main] coordinates {(7,0.2) (7,-0.2)};
    \node[main, above] at (axis cs:5,0) {\tiny\sffamily CN};
    \node[main, above] at (axis cs:7,0) {\tiny\sffamily CN};
    \draw[main, decorate, decoration={brace, raise=0.2cm, mirror}] (5,0) -- (7,0);
    \node[main, below] at (axis cs:6,-0.6) {\tiny\sffamily all CNs};
    % absolute max
    \addplot[main] coordinates {(9,0.2) (9,-0.2)};
    \node[main, below] at (axis cs:9,-0.2) {\tiny\sffamily CN};
    \draw[main, ultra thick] (5,0) -- (7,0);
    \fill[supp] (axis cs:9,4) circle[radius=0.2];
    \addplot[supp, dotted] coordinates { (9,4) (0,4) } node[above right] {\tiny\sffamily its \(y\)-value is the abs max};
    \addplot[warn] coordinates {(10,0.2) (10,-0.2)};
    \node[warn, above] at (axis cs:10,0.2) {\tiny\sffamily \xcancel{CN}};
  \end{axis}
\end{tikzpicture}
\end{document}
